\documentclass[10pt]{article}
\usepackage[margin=2.6cm]{geometry}
\usepackage{amsmath,amssymb,amsthm}
\usepackage{booktabs}
\usepackage[hidelinks]{hyperref}

\theoremstyle{definition}
\newtheorem{definition}{Definition}
\newtheorem{lemma}{Lemma}
\newtheorem{conjecture}{Conjecture}
\newtheorem*{remark}{Remark}

\title{\bfseries The universal rung profile $\Delta(\ell)$:\\
a lemma candidate}
\author{Denis Joubert\\ \small Independent researcher, Switzerland
\and Grok\\ \small draft for the author, 1 September 2026}
\date{September 1, 2026 --- lemma campaign}

\newcommand{\R}{\mathbb{R}}
\begin{document}
\maketitle

\begin{abstract}
\noindent
The quasi-null ladders of the fifteen truncated Weil forms share a spacing
profile $\Delta(\ell)$ to $\pm12\%$. The dilated-Slepian model for that
profile has no well-defined exponent and is not the lemma. What the
thirty-three rungs actually say is sharper and weaker: after a short
rise, spacings \emph{saturate} at $\Delta_\infty\approx16$ nats,
independently of the $L$-function. The $\sqrt{\ell}$ fit is a two-parameter
compromise with a ramp-then-plateau, not an asymptotic. We state the
universality statement as a lemma candidate in pure harmonic analysis
(approximation numbers of windowed Fourier restriction, after scaling out
$\lambda_0$), record the factorization into one arithmetic speed $s(\chi)$
times one window-invariant profile, and isolate the two assertions a
proof must supply: (i) scale invariance of the deflated problem past the
infrared turning point; (ii) character-blindness of the leftover
operator. No claim is made that the lemma is proved.
\end{abstract}

\section{The object}

Fix $\mu=e^L$ and a real primitive character $\chi$ (or $\zeta$). Let
$Q_{\chi,\mu}$ be the semi-local Weil form on even test functions
supported in $I_L=[-L/2,L/2]$, assembled as in~\cite{Quorum}. Write
$0<\lambda_0\le\lambda_1\le\cdots$ for its eigenvalues and
\[
\ell_k \;=\; -\ln\lambda_k,\qquad
\Delta_k \;=\; \ell_{k+1}-\ell_k \;=\; \ln\frac{\lambda_k}{\lambda_{k+1}}.
\]
Call $\ell_k$ the \emph{depth} of rung $k$ and $\Delta_k$ its
\emph{spacing} (nats). The ground-state law
$\ell_0=s(\chi)\,\mu+b(\chi)$ is \emph{not} part of this note: $s(\chi)$
is arithmetic. The question here is the internal architecture of one
spectrum, at fixed window.

Across $33$ rungs of fifteen $L$-functions (characters at $\mu=38$,
$\zeta$ at $\mu=16$ in a depth-adequate basis) the spacings collapse by
level bands~\cite{Depth,LabNB}:
\begin{center}\small
\begin{tabular}{lcc}
\toprule
level band $\ell$ & $\Delta$ & sample \\
\midrule
$[0,15)$  & $10.9\pm1.5$ & eight rungs, eight characters \\
$[15,40)$ & $13.5\pm1.5$ & \\
$[40,80)$ & $16.0\pm1.6$ & $\zeta$ interleaves with $\chi_3,\chi_4$ \\
$[80,150)$& $16.0\pm2.1$ & \\
\bottomrule
\end{tabular}
\end{center}
A two-parameter fit $\Delta(\ell)\approx9.8+0.65\sqrt{\ell}$ has RMS
$1.8$ and inter-character scatter $\sim12\%$. A dilated-Slepian reading
$\lambda_k\approx(1-\lambda_k^{\mathrm{Slepian}})^\kappa$ produced
$\kappa=2.85\pm0.26$ at constrained $c$, then $3.46$ / $4.48$ / $5.37$
on ten-rung ladders at free $c$~\cite{LabNB}: the model has no
well-defined exponent.

\section{What the bands say}

The four bands are not a slow square-root climb. They are a
\emph{ramp of $\sim40$ nats} followed by a \emph{plateau}:
\[
\Delta(\ell) \;\longrightarrow\; \Delta_\infty \;\approx\; 16
\qquad(\ell\gtrsim40).
\]
On the plateau each successive eigenvalue is larger than the previous
one by a constant factor $e^{16}\approx9\times10^6$. That is a geometric
tail, not a Slepian tail and not a $\sqrt{\ell}$ tail.

The $\sqrt{\ell}$ fit is what a two-parameter curve does when it is
asked to interpolate a step of height $5$ over a run of length $40$.
It is a useful mnemonic. It is not the lemma.

\begin{remark}
A geometric tail $\ell_k=\ell_0+k\Delta_\infty$ is the generic
signature of a \emph{scale-invariant leftover problem}: after the
infrared desert and the turning-point region have been spent on the
first rungs, each new orthogonality constraint multiplies the Rayleigh
quotient by a constant. The constant is an invariant of the window,
not of the zero set.
\end{remark}

\section{The lemma candidate}

\begin{definition}[Windowed restriction]
Let $V_L$ be the space of even real functions supported in $I_L$,
with the $L^2$ inner product. For a discrete even set
$\Gamma\subset\R$ write
$Q_\Gamma(f)=\sum_{\gamma\in\Gamma}|\hat f(\gamma)|^2$.
When $\Gamma$ is the nontrivial zeros of $L(s,\chi)$ (and the pole
term, for $\zeta$), $Q_\Gamma$ on $V_L$ \emph{is} the truncated Weil
form, by the exact pairing of~\cite{Quorum}.
\end{definition}

\begin{conjecture}[Universality of spacings]\label{conj:univ}
There exists a function $\Delta_\infty:[0,\infty)\to(0,\infty)$,
depending on the window construction $(V_L,\widehat{\,\cdot\,})$ and
not on $\chi$, such that for every real primitive $\chi$ and every
depth-adequate cutoff $\mu$,
\[
\Delta_k(\chi,\mu)
\;=\;
\Delta_\infty(\ell_k(\chi,\mu)) \;+\; \varepsilon_{k,\chi,\mu},
\]
with $\mathrm{rms}(\varepsilon)/\Delta \lesssim 12\%$ on the thirty-three
rungs already measured, and $\Delta_\infty(\ell)\to\Delta_\infty\approx16$
as $\ell\to\infty$.
\end{conjecture}

This is the surviving content of~\cite{Depth} \S4, stripped of the
dissolved $\kappa$-model. It is a statement in harmonic analysis:
approximation numbers of a Fourier restriction to a discrete set,
\emph{after the first eigenvalue has been scaled out}. The arithmetic
of $\chi$ is allowed to set $\lambda_0$. It is not allowed to set the
ratios $\lambda_k/\lambda_{k+1}$ once depth is given.

\begin{lemma}[Factorization --- formal consequence]\label{lem:fact}
Assume Conjecture~\ref{conj:univ} and the affine ground-state law
$\ell_0=s(\chi)\mu+b(\chi)$. Then the whole quasi-null ladder is
determined by two arithmetic scalars $(s,b)$ and one window-invariant
profile $\Delta_\infty$:
\[
\ell_0=s\mu+b,\qquad
\ell_{k+1}=\ell_k+\Delta_\infty(\ell_k).
\]
Equivalently: $\chi$ chooses a starting depth; the window chooses the
staircase.
\end{lemma}

Lemma~\ref{lem:fact} is not an additional claim. It is the
factorization already measured: architecture $\times$ speed. The
content that still needs a proof is Conjecture~\ref{conj:univ}.

\section{What a proof must supply}

Two harmonic-analysis assertions, neither mentioning RH.

\paragraph{A1. Scale invariance past the turning point.}
Let $R_k$ be the Rayleigh problem for $Q_\Gamma$ on the orthogonal of
the first $k$ eigenfunctions in $V_L$. For $\ell_k\gtrsim40$, $R_k$ is
equivalent, up to a scalar, to $R_{k+1}$ after a universal
renormalization of $V_L$. In particular the minimax value drops by a
factor $e^{-\Delta_\infty}$ independent of $k$ and of $\Gamma$.

Heuristic: the first rungs spend the infrared desert (a Slepian /
Airy / Landau--Widom turning point, character-dependent through
$\gamma_1$). What remains is a problem on the \emph{window} ---
time-limited functions, Fourier mass expelled to the band frontier
(hyper-nullity~\cite{LabNB} \S16.3) --- whose next mode costs a
constant factor. The plateau $\Delta_\infty\approx16$ is that factor
in nats.

\paragraph{A2. Character-blindness of the leftover.}
Once $\lambda_0$ is scaled out, the leftover operator on
$v_0^\perp\subset V_L$ depends on $\Gamma$ only through quantities
already absorbed in $\lambda_0$. Distinct zero sets with the same
window produce the same ratios $\lambda_k/\lambda_{k+1}$ at equal
depth.

Heuristic: hyper-nullity says the residual mass of a deep rung does
not live on the interior zeros. If the leak is a frontier phenomenon,
the interior arithmetic of $\Gamma$ cannot set the next ratio.

A1 + A2 $\Rightarrow$ Conjecture~\ref{conj:univ}. Recruitment (rung
$k$ seats the $(k{+}1)$-st $\chi$-supported prime) is a statement
about \emph{organization}, not about \emph{spacing}, and is not
needed for the lemma. Mixing the two was the $\kappa$-trap.

\section{What the lemma is not}

\begin{itemize}
\item \emph{Not} $\lambda_k=(1-\lambda_k^{\mathrm{Slepian}})^\kappa$.
Classical prolate spacings in $-\ln(1-\lambda)$ units are $3$--$5$;
Weil spacings are $11$--$19$. A free $\kappa$ fits any one ladder and
does not transfer~\cite{LabNB} \S14.8.
\item \emph{Not} Landau--Widom for $\lambda_{\min}(L)$ as the window
grows~\cite{Chuk}. That law lives on the other axis (depth of the
ground state versus $L$). The $s\mu$ law already occupies that axis
and is linear in $\mu=e^L$, i.e.\ exponential in the window width ---
the lattice signature. $\Delta(\ell)$ is the internal mesh of one
matrix.
\item \emph{Not} a function of the prime gaps. Log-gaps of primes
shrink; $\Delta$ saturates. Recruitment assigns primes to rungs;
spacings are not those gaps.
\item \emph{Not} a counting functional of the zeros. Vacancy is
quasi-universal~\cite{Depth}; depth is a bandwidth phenomenon. The
lemma lives on that side.
\end{itemize}

\section{A numerical lock for $\Delta_\infty$}

The plateau rests on two bands and a handful of deep rungs ($\zeta$,
$\chi_3$, $\chi_4$). Three cheap checks would freeze or kill it.

\begin{enumerate}
\item Ten-rung ladders at one more deep character already computed
($\chi_5$ at $\mu=38$): if the last four spacings sit at $16\pm2$,
the plateau is a family fact, not a $\zeta$--$\chi_3$ coincidence.
\item A synthetic $Q_\Gamma$ with $\Gamma$ a rigid unfolding of
spacing $2\pi$ (no arithmetic, no desert variation): the same
$\Delta_\infty$ would confirm A2 directly.
\item Dependence on $L$ at fixed $\chi$. If $\Delta_\infty$ is a
window invariant it must be stable when $\mu$ changes and only
$s\mu$ slides the starting depth. A counter-variation with $L$
would mean $\Delta_\infty$ still carries a piece of $s(\chi)$.
\end{enumerate}

Until (1)--(3) are run, $\Delta_\infty\approx16$ is a reading of four
bins, not a constant.

\section{Status}

Conjecture~\ref{conj:univ} is the lemma candidate. Lemma~\ref{lem:fact}
is bookkeeping. A1 and A2 are the proof obligations. The $\sqrt{\ell}$
fit and the $\kappa$-model are retired as statements; they remain
useful as the history of a failed identification.

Nothing here is claimed about RH. The lemma, if true, says that once
the ground state has been drilled to depth $s\mu$, the rest of the
quasi-null ladder is a property of the window.

\begin{thebibliography}{9}\small
\bibitem{Depth} D.~Joubert, \emph{Depth phenomenology of truncated Weil
quadratic forms: fifteen $L$-functions, one ladder}, same repository
(2026).
\bibitem{Quorum} D.~Joubert, \emph{Proper sub-Euler products violate
Weil positivity on a fixed window: certified witnesses}, same
repository (2026).
\bibitem{LabNB} D.~Joubert, \emph{Le milieu des nombres premiers}, lab
notebook, same repository (2026), \S14.2--14.8 and \S16.3.
\bibitem{Chuk} M.~Chuk, \emph{Weil positivity in compact windows:
certified two-sided bounds and a Landau--Widom decay law},
arXiv:2608.24827 (2026).
\end{thebibliography}

\section*{Addendum --- leakage profile, 1 September afternoon}

Notebook \S16.4--16.7. Rungs 0--2 share one leakage shape, scaled by
depth: evidence for A1 on the zero side. D2 (a $\mu$-independent
$\tau$) died at the 85-digit zeros; A1 is scale invariance of the
\emph{deflated} problem, not constancy of $\tau$. The leftover still
carries arithmetic through $s(\chi)$ via the endpoint law
$|\hat v_0(\gamma_1)|\approx C\lambda_0$. A2 remains open.
\end{document}
